% =============================================================
% Introduction générale — Mémoire ORACLE
% Version compacte (cible ~3 pages)
% =============================================================

\documentclass[14pt,a4paper,openany]{extreport}

\usepackage[T1]{fontenc}
\usepackage[utf8]{inputenc}
\usepackage[french]{babel}
\addto\captionsfrench{%
  \renewcommand{\tablename}{Tableau}%
  \renewcommand{\figurename}{Figure}%
}

\usepackage{lmodern}
\usepackage{setspace}
\onehalfspacing
\usepackage[a4paper,inner=3cm,outer=2.3cm,top=2.5cm,bottom=2.7cm,headheight=15pt,bindingoffset=0.3cm]{geometry}

\usepackage{amsmath,amssymb,amsthm}
\usepackage{graphicx}
\usepackage{xcolor}
\definecolor{uridarkblue}{RGB}{30,40,90}
\definecolor{urigray}{RGB}{70,70,70}
\definecolor{uriaccent}{RGB}{180,60,40}

\usepackage{array}
\usepackage{booktabs}
\usepackage{tabularx}
\usepackage[most]{tcolorbox}

\usepackage{titlesec}
\titleformat{\section}
  {\normalfont\Large\bfseries\color{uridarkblue}}
  {\thesection.}{0.8em}{}

\usepackage{fancyhdr}
\pagestyle{fancy}
\fancyhf{}
\fancyhead[L]{\textsc{\nouppercase{\leftmark}}}
\fancyhead[R]{\thepage}
\fancyfoot[L]{\footnotesize Mémoire de Master en Informatique, Université de Dschang}
\fancyfoot[R]{\footnotesize URIFIA}
\renewcommand{\headrulewidth}{0.4pt}
\renewcommand{\footrulewidth}{0.2pt}

\usepackage[font=small,labelfont={sc,bf},textfont=it]{caption}
\usepackage[numbers,square,sort&compress]{natbib}
\usepackage[hidelinks]{hyperref}

\newcommand{\Oracle}{\textsc{Oracle}}
\newcommand{\CorrDiff}{\textsc{CorrDiff}}
\newcommand{\HR}{\ensuremath{\text{HR}}}
\newcommand{\LR}{\ensuremath{\text{LR}}}
\newcommand{\muHR}{\ensuremath{\mu_{\HR}}}
\newcommand{\Adag}{\ensuremath{A_{\mathrm{dag}}}}

\setcounter{page}{1}
\pagenumbering{roman}

\begin{document}

\chapter*{Introduction générale}
\addcontentsline{toc}{chapter}{Introduction générale}
\markboth{Introduction générale}{Introduction générale}

\section*{Contexte et motivation}

Le changement climatique impose des décisions à des échelles spatiales que les modèles climatiques globaux ne fournissent pas immédiatement. Ces modèles simulent l'atmosphère à une résolution de l'ordre de la centaine de kilomètres, suffisante pour les grands modes de variabilité, mais très en deçà de ce qu'exige la planification locale d'infrastructures, l'agriculture pluviale ou la gestion du risque hydrologique. La précipitation présente une structure spatiale qui se déploie sur quelques kilomètres et qui interagit fortement avec le relief : un modèle qui produit une valeur moyenne sur cent kilomètres ne peut pas restituer cette structure. C'est cet écart d'échelle que cherche à combler le \emph{downscaling} climatique.

La motivation finale est l'aide à la décision dans des régions à forte vulnérabilité hydrologique et faible densité d'observations, dont l'Afrique subsaharienne. Construire et valider une telle architecture directement sur le continent bute aujourd'hui sur deux verrous croisés : l'Afrique héberge moins de $1\%$ des capacités HPC mondiales~\cite{bachmann_hpc_africa_2025}, ce qui exclut la production locale de RCM denses, et il n'existe pas à ce jour de jeu de référence HR suffisamment dense pour entraîner un downscaler génératif causal sur la zone visée.

Faute de pouvoir valider la méthode \emph{in situ}, nous l'éprouvons sur un terrain qui en éprouve simultanément les trois axes du trilemme : la Nouvelle-Zélande, choisie par Rampal et~al.~\cite{rampal_cgan_2024} pour le même argument et que nous traitons ici comme \emph{bac à sable exigeant}. Le pays concentre, sur une emprise de moins de $270\,000$~km², plusieurs régimes climatiques (subtropical au nord, marin tempéré à l'ouest, semi-aride sous le vent à l'est, alpin en altitude), et la chaîne des Southern Alps engendre un des gradients orographiques de précipitation les plus extrêmes au monde --- conditions qui sollicitent en même temps la stabilité d'entraînement, le réalisme des extrêmes et la généralisation spatiale. Une architecture qui résiste à l'ablation causale dans ce cadre fournit une preuve forte d'\emph{universalité architecturale}, indépendante du climat d'entraînement ; l'application opérationnelle à l'Afrique reste l'objectif final et est portée comme axe~5 des perspectives, conditionnée à la disponibilité d'un jeu HR de référence sur la zone visée. Les outils existants ne fournissent pas ce service de manière satisfaisante : ils sont soit déterministes (donc impropres au raisonnement sur le risque), soit insuffisamment résolus, soit non transférables hors de leur domaine de calibration~\cite{maraun_widmann_2018}.

\section*{Problématique}

Les architectures profondes récentes (GAN résiduels, diffusion à deux étages) ont amélioré la qualité prédictive et la calibration probabiliste, mais restent fragiles sous changement de distribution. Un modèle entraîné sur le climat historique apprend des associations contingentes qu'aucune régularisation a posteriori ne garantit de redresser. La causalité structurelle est, en théorie, une réponse de principe : un modèle organisé autour de relations causales identifiées devrait mieux transférer en inter-GCM. Cette promesse n'a pas été opérationnalisée dans un downscaler génératif au-delà du diagnostic post-hoc.

Le paysage tient en un trilemme : aucune approche existante ne garantit simultanément la stabilité d'entraînement, le réalisme des extrêmes et la robustesse inter-GCM historique. Le paradigme \emph{two-stage} de Mardani et~al.~\cite{mardani_corrdiff_2023} adresse les deux premières propriétés, pas la troisième : son premier étage est un réseau de régression générique, aussi vulnérable au changement de climat qu'un modèle monolithique. Question centrale du mémoire : \emph{est-il possible de causaliser le premier étage du paradigme two-stage de manière à doter le downscaling probabiliste d'une robustesse de principe, sans dégrader la stabilité ni le réalisme acquis} ?

\section*{Questions de recherche, objectifs et hypothèses}

Trois questions structurent la suite. \textbf{Q1.} Les architectures génératives existantes (cGAN résiduel, \CorrDiff{}) suffisent-elles à garantir une robustesse inter-GCM en régime historique ? \textbf{Q2.} Un modèle causal structurel intégré dans une dynamique récurrente permet-il de mieux guider le downscaling probabiliste, et selon quelle interface avec un décodeur de diffusion résiduelle ? \textbf{Q3.} Quelles métriques permettent de discriminer une véritable causalité architecturale d'une simple structuration corrélationnelle ?

La conception du modèle obéit à six objectifs notés O1 à O6, repris au chapitre III : \textbf{O1} stabilité d'entraînement, \textbf{O2} réalisme des extrêmes, \textbf{O3} causalité architecturale (la matrice $\Adag$ doit être indispensable à la prédiction par construction), \textbf{O4} robustesse inter-GCM en régime historique, \textbf{O5} efficacité computationnelle, \textbf{O6} interprétabilité.

Trois hypothèses de travail sont formulées en début de projet. \textbf{H1.} Un modèle apprenant des relations structurelles sous contrainte d'acyclicité, dans une dynamique récurrente, généralise mieux qu'un modèle purement corrélatif. \textbf{H2.} La séparation entre moyenne causale déterministe et résidu stochastique stabilise l'entraînement et améliore la calibration. \textbf{H3.} L'indispensabilité architecturale du DAG limite le risque qu'il ne soit qu'un décor inopérant, sans pour autant garantir une identification causale au sens strict.

\noindent\textbf{Précision terminologique.} L'expression « causalité architecturale » désigne une \emph{régularisation structurelle inductive} : la matrice $\Adag$ est indispensable au chemin de calcul (ablation $\Rightarrow$ chute de performance), ce qui prouve une dépendance fonctionnelle, non une identification causale au sens de Pearl~\cite{pearl_causality_2009}. La contrainte d'acyclicité DAGMA est un dispositif de traitabilité ; la physique atmosphérique réelle comporte des boucles de rétroaction qu'un graphe acyclique n'encode pas. Le DAG appris doit donc être lu comme un \emph{graphe de contraintes fonctionnelles}.

H1 sera évaluée par comparaison CRPS et signes interventionnels, H2 par les ratios spread/skill et la convergence d'entraînement, H3 par l'ablation de la matrice apprise. Le bilan détaillé figure en conclusion.

\section*{Contributions et organisation}

Le mémoire propose quatre contributions. \textbf{C1.} Une architecture \emph{two-stage} causale, \Oracle{}, qui place la causalité architecturale dans un chemin de calcul indispensable plutôt que dans une régularisation de la perte. \textbf{C2.} Un réseau causal récurrent (RCN) intégrant dans une seule cellule un modèle causal structurel dynamique, une mise à jour récurrente GRU et une pénalité d'acyclicité DAGMA~\cite{bello_dagma_2022}. \textbf{C3.} Une interface inter-étages par concaténation de canaux, qui garantit qu'aucune information ne contourne le chemin causal. \textbf{C4.} Un protocole expérimental centré sur la calibration probabiliste, les extrêmes et les ablations causales --- incluant un test d'intervention $Q_{\mathrm{int}}$ sur des variables atmosphériques, original dans le contexte du downscaling profond.

Le manuscrit comporte quatre chapitres et une conclusion. Le \textbf{chapitre I} pose les fondements du downscaling probabiliste et formule le trilemme. Le \textbf{chapitre II} parcourt l'état de l'art par familles (déterministes, GAN, vraisemblance, diffusion) et identifie le \emph{gap} scientifique. Le \textbf{chapitre III} formalise \Oracle{} : graphe atmosphérique hétérogène, encodeur, RCN, apprentissage du DAG sous DAGMA, décodeur graphe-à-grille, interface inter-étages, diffusion résiduelle EDM. Le \textbf{chapitre IV} couvre les matériels, méthodes, expérimentations et discussion. La \textbf{conclusion} synthétise les contributions, répond aux questions de recherche et trace les perspectives.

\end{document}
